% GENERATED_BY: run_oc_core_monograph_translation_rework_dispatch_v1.py
% AI_ASSISTED_TRANSLATION_DRAFT: TRUE
% THEOREM_REVIEW_STATUS: APPROVED
% THEOREM_REVIEW_MODE: FORMAL_AI_GATE__CODEX_CLI_CHATGPT
% THEOREM_REVIEW_REVIEWER_ID: CODEX_CLI_CHATGPT__AUTO_DISPATCH_20260406
% THEOREM_REVIEW_AUTHORITY_CLASS: FINAL_ADVERSARIAL_EXTERNAL_LLM_REVIEWER
% THEOREM_REVIEWED_AT: 2026-04-14T12:38:36Z
% SOURCE_LANGUAGE: EN
% TARGET_LANGUAGE: DE
% SOURCE_EN_REF: content/frontmatter_oc_core_1_3_master.tex
\begin{titlepage}
\thispagestyle{empty}
\begin{center}
{\Large Ontology of Continua}\\[0.6cm]
{\huge \textbf{Core 1.3}}\\[0.3cm]
{\Large Quellenpriorisierte Master-Monographie}\\[1.2cm]

{\Large Alexander Yashin}\\[0.2cm]
{\normalsize Unabhängiger Forscher, Leipzig/Halle, Deutschland}\\[0.2cm]
{\normalsize \href{https://orcid.org/0009-0008-6166-0914}{ORCID 0009-0008-6166-0914}}\\[1.0cm]

{\normalsize Version 1.3}\\[0.2cm]
{\normalsize 13. April 2026}\\[1.2cm]

\begin{minipage}{0.88\textwidth}
\small
\textbf{Publikationsrolle.} Dieser Band ist die kanonische Master-Monographie für OC Core 1.3.
Er stellt die vollständige wissenschaftliche Hülle im Monographiemaßstab wieder her, behält den geerbten formalen Korpus von
Core 1.2 als quellengebundenes Substrat bei und ergänzt die für die
Freigabe 1.3 erforderlichen expliziten Ebenen für Quellenaudit, Theorem-Verbleib,
Chronologie, Theorem-Fahrplan, ausgearbeitete Beispiele und Publikationsgrenzen.

\medskip

\textbf{Versprechen an die Leserschaft.} Das Ziel dieses Dokuments ist nicht, das Rahmenwerk in ein
kurzes Paket zu komprimieren. Das Ziel ist, der Leserschaft einen inspizierbaren wissenschaftlichen Band zu geben, in dem
Definitionen, Theoremaussagen, Beweise, Abbildungen, Begrenzungen und Provenienz in
ihrer natürlichen Reihenfolge gelesen werden können, ohne die Theorie aus Übersichtsartefakten
oder verborgenen Begleitregistern rekonstruieren zu müssen. Der Band enthält daher einen
expliziten Leserleitfaden, einen Theorem-Fahrplan, ausgearbeitete Beispiele und einen Anhang zur
Gutachternavigation, sodass ein feindseliger Gutachter und eine ernsthafte Nichtfachperson
über unterschiedliche rechtmäßige Pfade in dasselbe wissenschaftliche Objekt eintreten können,
ohne die Beweisgrenze zu verändern.
\end{minipage}

\vfill

% DEDICATION_REQUIRED_DO_NOT_REMOVE
% OC_CORE_PUBLICATION_DEDICATION
{\small\itshape Meiner lieben Ehefrau Maria gewidmet, ohne die diese Arbeit nicht möglich gewesen wäre.}\\[0.8cm]

{\small Offener Grundlagenpfad: github.com/alexanderyashin/ontology-of-continua-core-main}\\
{\small Archivierter Vorgängernachweis: Zenodo DOI 10.5281/zenodo.17903912}
\end{center}
\end{titlepage}

\begin{abstract}
Core 1.3 ist die quellenpriorisierte, erneut validierte Master-Monographie der Ontology of Continua.
Sie geht von der vollständigen formalen Hülle und dem Kapitelbaum aus, die Core 1.2 zur ersten
axiomatisch abgeschlossenen OC-Freigabe machten, und unterzieht diesen geerbten Korpus anschließend einer strengeren
Publikationsdisziplin: expliziten Quellennachweisen, einer Trennung nach Theorem-Verbleib, chronologiebewussten Reparaturhinweisen,
Grenzen des Nicht-Behaupteten und einer für die externe Publikation geeigneten wissenschaftlichen Aufbereitung.

Das Ergebnis ist bewusst umfangreicher als der daraus extrahierte kurze, auf Zeitschriftenpublikation ausgerichtete Artikel.
Diese Monographie ist das didaktische und formale Masterdokument: Sie trägt die Strukturtheorie, das
axiomatische und theoremtragende Rückgrat, die Hierarchie der Kontinua und Einbettungsräume, die
ebenenübergreifenden und domänenspezifischen Module, die Falsifizierbarkeitsprogramme sowie die expliziten
1.3-Materialien, die erläutern, was unverändert fortbesteht, was revidiert wird, was begrenzt ist und was zurückgestellt wird.

Die wissenschaftliche Rolle des Bandes ist daher doppelt. Erstens bewahrt er die Lesbarkeit auf
Monographieebene, die ein ernstzunehmender Grundlagenrahmen erfordert. Zweitens liefert er die kanonische
Quelle, aus der engere zeitschriftenorientierte Kernartefakte und kanalspezifische Artefakte abgeleitet werden können,
ohne die verteidigte Behauptung stillschweigend zu erweitern oder zu verengen. Die Monographie ist damit
sowohl ein eigenständiges wissenschaftliches Werk als auch das autoritative Substrat für externe
Publikation, Begutachtung, Übersetzung und spätere Projektion in Produkt-, Geschäfts- und
öffentliche Kanäle.

Die Freigabe 1.3 macht zudem die Beweisnavigation explizit. Anstatt vorauszusetzen, dass jede Leserin und jeder Leser
denselben Pfad durch eine große formale Hülle finden wird, enthält der Band nun einen eigenen
Theorem-Fahrplan, ausgearbeitete Beispiele zur Semantik von Kollaps / Residuum / Wiedergeburt und zur Doktrin der K-Ebenen
sowie eine Matrix zur Gutachternavigation, die Einwände, Beweisziele und Quellennachweise in einen
inspizierbaren Pfad abbildet.
\end{abstract}

\clearpage
